\documentclass[12pt]{article}
\usepackage{sbc-template}
\usepackage[T1]{fontenc}
\usepackage[utf8]{inputenc}
\usepackage[brazil]{babel}
\usepackage{graphicx,url}
\sloppy

\title{MindScreen AI: uma abordagem exploratória para análise de tempo de tela e indicadores de saúde mental na infância e adolescência}
% Substitua os campos abaixo pelos dados reais antes da submissão.
\author{Nome do(a) autor(a)\inst{1}}
\address{Instituição e departamento -- Cidade -- Estado -- Brasil
  \email{email@instituicao.br}}

\begin{document}
\maketitle

\begin{abstract}
This paper presents MindScreen AI, a research prototype for studying associations between screen time and parent-reported current depression and anxiety in children and adolescents. The project harmonizes 2020--2024 National Survey of Children's Health data and implements exploratory analyses, predictive models, explanations and an interactive interface. The harmonized sample contains 166,343 records aged 6--17. Unweighted complete-case analyses suggest associations, but do not establish causality or population estimates. The prototype is not a diagnostic tool.
\end{abstract}

\begin{resumo}
Este artigo apresenta o MindScreen AI, protótipo de pesquisa sobre associações entre tempo de tela e relatos de responsáveis sobre depressão e ansiedade atuais em crianças e adolescentes. O projeto harmoniza dados da National Survey of Children's Health de 2020 a 2024 e implementa análises exploratórias, modelos preditivos, explicações e interface interativa. A amostra harmonizada contém 166.343 registros de 6 a 17 anos. Análises não ponderadas de casos completos sugerem associações, sem estabelecer causalidade ou estimativas populacionais. O protótipo não é uma ferramenta de diagnóstico.
\end{resumo}

\section{Introdução}
O uso de telas integra o cotidiano de crianças e adolescentes, e sua relação com indicadores de saúde mental requer análise cuidadosa de fatores contextuais e das limitações dos dados observacionais. A disponibilidade de microdados da \textit{National Survey of Children's Health} (NSCH), realizada nos Estados Unidos, permite investigar padrões em respostas fornecidas por responsáveis \cite{census2020,census2024}. Contudo, relações observadas em uma pesquisa transversal não demonstram que a exposição precedeu ou causou um desfecho.

O objetivo deste trabalho é apresentar o MindScreen AI, um protótipo que organiza microdados da NSCH, investiga associações entre categorias de tempo de tela e dois indicadores de saúde mental e disponibiliza recursos de modelagem e visualização. As tarefas de classificação de depressão e ansiedade atuais são independentes. As contribuições consistem na documentação de um processo de harmonização entre edições, em uma análise amostral ajustada e na implementação de uma infraestrutura reprodutível de experimentos. Nenhuma saída do sistema constitui diagnóstico médico.

\section{Dados e harmonização}
São utilizadas as edições \textit{topical} da NSCH de 2020 a 2024, com população analítica de 6 a 17 anos. A escolha busca um universo comum para as variáveis inicialmente consideradas. As listas e frequências oficiais de cada edição orientam as regras e verificações; as listas de 2020 e 2024 exemplificam a documentação consultada \cite{census2020,census2024}. O projeto preserva os arquivos de origem separados dos dados processados.

Cada arquivo é submetido a verificações técnicas de colunas, identificadores, domínios e frequências selecionadas, incluindo comparações de pesos e distribuições com publicações do Census. Um manifesto vincula o dataset produzido a seu hash e registra os preditores permitidos. Essas verificações não equivalem a revisão científica independente nem autenticam integralmente a origem dos arquivos.

A variável \texttt{SCREENTIME} é mantida em cinco categorias, de menos de uma hora a quatro horas ou mais em dias de semana. A documentação de 2024 explicita a exclusão de tarefas escolares, expressão ausente das listas anteriores consultadas; essa diferença documental não comprova, por si só, mudança no instrumento. Os preditores iniciais também incluem idade, sexo, sono, atividade física e dificuldade para fazer amizades. Variáveis de bullying e itens comportamentais são mantidos para exploração, mas excluídos do conjunto inicial de preditores por diferenças documentais ou proximidade com sintomas.

Os desfechos são construídos separadamente a partir das perguntas sobre histórico e condição atual de depressão ou ansiedade. Um histórico negativo com salto documentado é tratado como negativo; quando o histórico é positivo, a resposta sobre condição atual determina o desfecho. Respostas inconsistentes ou insuficientes permanecem indeterminadas. Os componentes dos desfechos são armazenados separadamente e não são utilizados como preditores. Ausências preservam os motivos disponíveis; campos vazios de CSV não permitem recuperar motivos que tenham sido perdidos na conversão.

\section{Métodos}
\subsection{Associação ajustada}
Para cada desfecho, uma regressão logística não penalizada é ajustada aos casos completos de 2020 a 2023. As categorias de tempo de tela são comparadas à categoria de menos de uma hora, com ajuste simultâneo por idade, sexo, sono, atividade física, dificuldade para fazer amizades e ano. As razões de chances (\textit{odds ratios}, ORs) são exponenciais dos coeficientes e descrevem associações condicionais na amostra. A análise não usa pesos amostrais e não apresenta intervalos de confiança que representem o desenho complexo da pesquisa. Covariáveis contemporâneas podem ser mediadoras ou consequências dos desfechos; o ajuste não elimina confundimento nem causalidade reversa.

\subsection{Modelagem e avaliação}
O sistema implementa regressão logística, Random Forest e XGBoost, além de um classificador de referência constante. Comparam-se dois conjuntos de preditores: experimento A, com tempo de tela, e experimento B, sem essa variável. A seleção de hiperparâmetros utiliza a área sob a curva ROC média em cinco partições estratificadas dos dados de 2020 a 2023. Imputação e codificação são ajustadas dentro de cada partição para reduzir vazamento de informação. Os pesos amostrais são preservados nos dados, mas não são aplicados ao treinamento.

A implementação prevê avaliação temporal em 2024 com métricas de discriminação e classificação e bootstrap pareado entre experimentos, sem incorporar o desenho amostral. Na revisão documentada do protocolo, o limiar é escolhido por F1 em predições fora da partição de desenvolvimento e a busca do XGBoost passa a incluir balanceamento de classes. Essa revisão ocorreu após observação de desempenho em 2024; portanto, esse ano deixou de ser um teste totalmente cego para a decisão de protocolo. É necessário um novo conjunto independente para avaliação confirmatória.

\subsection{Explicabilidade e interface}
O projeto implementa explicações SHAP para modelos selecionados, com verificações de aditividade e agrupamento posterior de colunas codificadas por variável original \cite{lundberg2017}. As contribuições descrevem o comportamento do modelo em relação a uma amostra de referência; não medem efeitos de uma intervenção. A interface Streamlit reúne visualizações e um simulador de pesquisa que usa preditores aprovados. Pontuações não calibradas e rótulos são apresentados sem interpretação clínica. O carregamento de modelos depende de artefatos locais e verificações de integridade.

\section{Resultados disponíveis}
O relatório de harmonização presente no projeto registra 166.343 observações de 6 a 17 anos, distribuídas conforme a Tabela~\ref{tab:amostra}. As contagens são amostrais e não ponderadas. Na amostra completa, o desfecho de depressão apresenta 11.193 positivos, 154.185 negativos e 965 indeterminados; o de ansiedade apresenta 26.051 positivos, 138.650 negativos e 1.642 indeterminados. Os desfechos refletem relatos dos responsáveis e não diagnósticos produzidos pelo sistema.

\begin{table}[ht]
\centering
\caption{Registros por edição e após restrição etária.}
\label{tab:amostra}
\begin{tabular}{lrr}
\hline
Ano & Total na edição & 6--17 anos \\
\hline
2020 & 42.777 & 30.636 \\
2021 & 50.892 & 30.173 \\
2022 & 54.103 & 34.362 \\
2023 & 55.162 & 33.638 \\
2024 & 51.375 & 37.534 \\
\hline
\end{tabular}
\end{table}

A análise ajustada disponível inclui 124.795 casos completos para depressão e 124.299 para ansiedade. Em relação a menos de uma hora, as ORs amostrais para quatro horas ou mais são, respectivamente, 1,79 e 1,65. Essas estimativas não ponderadas não representam razões de risco, efeitos causais ou inferência populacional, e não há intervalos de confiança que permitam avaliar sua precisão sob o desenho da NSCH. A análise excluiu, respectivamente, 4.014 e 4.510 registros de 2020--2023 por ausência no desfecho ou nas covariáveis.

Há registros locais de uma execução anterior de treinamento, cuja configuração não inclui a busca de balanceamento do XGBoost prevista na revisão. Naquele registro, o XGBoost foi selecionado por validação cruzada no experimento A para ambos os desfechos. Como o relatório de avaliação temporal e os artefatos completos da revisão não estão disponíveis nesta cópia, não são apresentados resultados de teste de 2024 nem se atribuem métricas dessa execução anterior ao protocolo revisado. A comparação A/B e as explicações dependem da geração e revisão dos respectivos relatórios.

\section{Discussão e limitações}
As associações encontradas são compatíveis com diferenças entre categorias de tempo de tela na amostra analisada, mas não permitem concluir que alterar o tempo de tela alteraria a ocorrência dos desfechos. Medidas autorreferidas por responsáveis, dados ausentes, seleção de casos completos e possíveis diferenças entre edições podem introduzir vieses. A ausência de pesos e de estimadores de variância apropriados impede generalização estatística à população. Além disso, os dados se referem aos Estados Unidos, e não houve validação clínica, calibração individual ou avaliação de transporte para o Brasil.

O manifesto e os hashes auxiliam a rastreabilidade, mas não substituem uma revisão metodológica humana. A comparação entre modelos com e sem tempo de tela pode investigar contribuição preditiva nas condições do experimento, sem atribuir causalidade. A revisão após observar 2024 reforça a necessidade de avaliação independente posterior. Limites de páginas, identificação de autoria e demais exigências devem ser adequados ao evento pretendido.

\section{Conclusão}
O MindScreen AI reúne harmonização, análise de associação, experimentos preditivos e interface para pesquisa exploratória com dados da NSCH. Os resultados disponíveis demonstram a viabilidade do fluxo e indicam associações amostrais que exigem interpretação cautelosa. Como trabalhos futuros, destacam-se revisão científica dos dados e preditores, aplicação de métodos adequados ao desenho complexo, avaliação em dados independentes e verificação de calibração e transportabilidade antes de qualquer uso além da pesquisa.

\bibliographystyle{sbc}
\bibliography{referencias}
\end{document}
